% !TEX root = ../main.tex
% =====================================================================
% 4. PROTOCOL -- page budget 0.5
%
% Short on purpose. Its real job is to state the two things a reviewer
% would otherwise catch us on: the selection effect (issue I7) and the
% uneven seed counts (issue I4), BEFORE any result is shown.
% =====================================================================

\section{Experimental Protocol}
\label{sec:protocol}

\paragraph{Datasets.}
\webtraffic{} \citep{early2024millet} is synthetic -- $1{,}000$ series of length $1{,}008$ over $10$
classes -- and its generator records which time steps it modified to create each class. It is
therefore the only dataset here on which \ndcg{} can be computed at all, and every claim about
explanation \emph{correctness} in this paper rests on it. We also use the $85$ \ucr{} datasets
\citep{dau2019ucr} for which \millet{} published per-dataset results, and run the remaining $43$ so
that no conclusion depends on a convenient subset.

\paragraph{One configuration for everything.}
The central methodological choice of this study is that the baselines are \emph{re-trained here}
rather than quoted. A published accuracy mixes the effect of the architecture with the effect of the
training recipe its authors used, and the two cannot be separated afterwards. Under one identical
configuration they can. Every model in this paper -- ours, \millet{}, FCN and ResNet
\citep{wang2017fcn} -- is trained with Adam \citep{kingma2015adam} at learning rate
$1.25\times10^{-3}$, weight decay $1.2\times10^{-4}$, at most $400$ epochs with patience $60$, batch
size $\min(\lfloor n_{\mathrm{train}}/10\rfloor, 16)$, and a stratified $80/20$ validation split
whenever a dataset has at least $100$ training series. The full list is in \Cref{app:hparams}. The
price of this choice is that our numbers are not directly comparable with \millet{}'s published
ones; that comparison is a separate question, answered on its own in \Cref{app:published}, and
\Cref{sec:aopcr} shows why keeping the two apart is not merely tidiness.

\paragraph{Two things to hold against every result that follows.}
First, \textbf{\webtraffic{} was the screening set}. Each new encoder and head was evaluated there
before the expensive \ucr{} runs, and only the strongest went forward. The \webtraffic{} numbers in
\Cref{sec:cheap} therefore carry a selection effect, and the \ucr{} table in the same section is the
honest counterweight rather than an afterthought. Second, \textbf{seed coverage is uneven}: four
models were trained with three seeds and the other $68$ with one. Where three seeds exist the spread
is not negligible -- accuracy standard deviation ranges from $0.003$ to $0.030$, and AOPCR is
proportionally noisier still ($\pm 0.884$ for the baseline). We therefore adopt one rule and apply it
in both directions throughout: \emph{a difference smaller than the relevant seed spread is reported
as noise, not as a result}, and every single-seed number is marked as such.

\paragraph{Reporting.}
Classification is measured by accuracy, explanation quality by AOPCR and \ndcg{}, cost by parameter
count, model size, FLOPs and time per prediction. Over the \ucr{} archive we report the mean accuracy
\emph{and} the mean accuracy rank across datasets, following \citet{demsar2006statistical}, because a
mean over datasets this heterogeneous can be dominated by a few of them while a rank cannot; ranks
are computed over the $84$ datasets shared by all fully-swept models.
